% ---
% §7.3 Por qué la codificación cuántica no aumentó la separabilidad.
% Entradas que lo respaldan: D-029, D-035, H-013, H-031, H-036, H-037, H-041, H-043, H-044,
% H-047, H-048, H-053, H-054, H-055, H-057, D-044.
% ---

\subsection{Por qué la codificación cuántica no aumentó la separabilidad}
\label{sec:causas_separabilidad}

Los resultados apuntan a tres causas, que actúan en etapas distintas del \textit{pipeline}.

% fuente: H-031, H-037, H-039, D-029; 5b_concentracion_kernel.csv, 6_dimension_efectiva.csv
% respaldo: «The small geometric difference is a consequence of the exponentially large Hilbert space employed by existing quantum models, where all inputs are too far apart.» (huang2021power, p. 2)
\textbf{La codificación coloca a las lesiones casi ortogonales entre sí.} Con 12 qubits y ángulos en
$[0,\pi]$, la fidelidad mediana entre dos lesiones distintas fue de 0.002 a 0.003 en C4 y del orden
de $3 \times 10^{-5}$ en C5, frente a 0.61 del kernel RBF, y la dimensión efectiva de los kernels
cuánticos ocupó entre 168 y 188 de las 200 posibles. Cada lesión se parece a sí misma y a casi ninguna
otra: es el régimen que Huang et al. \cite{huang2021power} identifican como causa de una diferencia
geométrica pequeña, en el que todas las entradas quedan demasiado lejos entre sí dentro de un espacio
de Hilbert exponencialmente grande. Reducir la escala angular disipa la concentración, pero solo cuando
el término de interacción $x_i x_j$ deja de influir: a $c = 0.02$ el kernel cuántico es prácticamente
idéntico al de un mapa sin ese término, y el alineamiento que recupera es el que un kernel clásico ya
veía. Con 8 qubits la concentración disminuye y la conclusión se mantiene.

% fuente: H-044, H-053, H-054; X4_test.csv, X4_test_diferencias.csv, X7_resumen.csv, X8_test.csv
\textbf{La lectura de 12 valores esperados pierde información que el kernel conserva.} Un clasificador
de vectores de soporte sobre el kernel de C4 alcanzó en el conjunto de prueba 0.628 en masas y 0.759 en
calcificaciones, estadísticamente indistinguible del SVM con kernel RBF (0.649 y 0.769) y entre 0.10 y
0.13 por encima del perceptrón sobre los valores esperados $\langle Z_i\rangle$ de la misma condición.
El estado codificado retiene, por tanto, una información diagnóstica comparable a la clásica, y la
pérdida del Módulo 7 ocurre al reducir 4\,096 amplitudes a 12 observables locales. Esa pérdida no se
explica por el sorteo de los parámetros de la capa variacional, que en diez sorteos dejó a C4 y C5
entre 0.03 y 0.18 por debajo de C1, ni por no entrenarlos: entrenarlos junto con el perceptrón ganó a lo
sumo 0.05 de AUC, sin significancia.

% fuente: H-047, H-048; X6_anidada.csv, X6_etiquetas_construidas.csv
% respaldo: «When g is large, there is a potential for improvement over classical ML.» (huang2021power, p. 7)
% respaldo: «The engineered dataset is constructed via a simple eigenvalue problem with the exact procedure described in Supplementary Section 7 and the results are shown in Fig. 4.» (huang2021power, p. 8)
\textbf{Las etiquetas clínicas no se alinean con la geometría de los kernels cuánticos evaluados.} Una
búsqueda anidada con el mismo protocolo para ambos lados, sesgada a favor de lo cuántico en el número de
candidatos, no encontró ningún kernel cuántico que superara a los clásicos con las etiquetas reales.
Con un etiquetado diseñado según el procedimiento de Huang et al. \cite{huang2021power}, que asigna a las
mismas lesiones las clases que mejor se ajustan a la geometría del kernel cuántico y peor a la del clásico,
el kernel proyectado alcanzó en validación cruzada un AUC de 0.998 en masas y 0.994 en calcificaciones frente a
0.800 y 0.824 del mejor clásico. El protocolo, por tanto, sí detecta una ventaja cuando el etiquetado la
permite; su ausencia es una propiedad de estas etiquetas y de las familias de kernels probadas, no del
método de evaluación. El etiquetado diseñado no es, sin embargo, una vía para mejorar el diagnóstico: sus
clases se derivan del propio kernel cuántico y no de la patología confirmada, de modo que cambia el problema
en lugar de resolverlo. Su papel es el de un control positivo del protocolo; en un problema clínico la etiqueta
está fijada, y lo que puede diseñarse es la codificación (sección~\ref{sec:futuro_codificacion}).

% fuente: H-057, D-044; X11_transicion.csv, X11_umbral.csv, X11_etiquetas_disenadas.csv
La distancia entre ambos escenarios se puede medir. Las clases diseñadas coinciden con la patología
confirmada en el 50 al 56\,\% de las lesiones, como dos etiquetados no relacionados. Al sustituir de forma
gradual las etiquetas diagnósticas por las diseñadas, la ventaja cuántica aparece cuando entre el 40 y el
70\,\% de las lesiones siguen la geometría del kernel cuántico, según el kernel, y su tamaño crece con la
diferencia geométrica: el kernel proyectado llega a ganar 0.14 de AUC, y el de C4, como mucho 0.035. El
diagnóstico benigno o maligno, descrito con las 12 características radiómicas, está en el extremo
opuesto: en él, ninguna fracción de las etiquetas sigue esa geometría.

% fuente: H-041, H-043, H-053, H-055; X3_comparacion_k8_k12.csv, X9_resumen_familia_a.csv
Ninguna de las elecciones de diseño que podían haber producido el resultado lo explica por sí sola: la
conclusión se mantuvo con 8 qubits, con dos repeticiones de la capa variacional, con diez sorteos de sus
parámetros, con esos parámetros entrenados, con otras submuestras y con todo el conjunto de entrenamiento.
